\chapter{Design Guidance for Implementers}
\label{ch:implications}

A useful design starts with the quantity the user needs and works backward to the observations required to estimate it.
Ball launch, club delivery and simulated carry share data, but they are not interchangeable validation targets.
This chapter connects the sensing geometry of the earlier chapters, the implementation conditions of \cref{app:impl}, and the evidence standard of \cref{ch:accuracy}.
The resulting capability tiers are proposed development stages with acceptance conditions, not claims that a parts list already delivers commercial performance.

A possible starting architecture combines a continuous-wave Doppler module, one or more angle-sensitive radar modules, a synchronized impact trigger and an embedded host.
The OPS243-A and K-LD7 are examples discussed elsewhere in the review, not interchangeable building blocks: their observables, latency, sampling and access to raw data differ.
Demonstrating a valid radial-speed channel is a useful first milestone.
Calling this arrangement the cheapest accurate monitor, or equivalent to a commercial radar architecture, would require evidence about costs and performance that the component list does not supply.

\section{Ball Data: Match Each Claim to Its Observations}

\begin{enumerate}
\item \textbf{Ball speed.}
A Doppler observation constrains radial velocity, not automatically total ball speed.
Separate ball returns from club, shaft and multipath returns, and establish calibration, aliasing limits and timing before claiming an accuracy tolerance.
For a correctly associated target with radial velocity $v_r=V\cos\beta$, where $\beta$ is the angle between velocity and the sensor's line of sight, the local correction is
\begin{equation}
\label{eq:designspeeduncertainty}
V=v_r\sec\beta,\qquad
\delta V\approx\sec\beta\,\delta v_r+V\tan\beta\,\delta\beta.
\end{equation}
Use radians for the angle perturbation and retain covariance between the inputs when propagating uncertainty.
The correction becomes ill-conditioned near perpendicular motion.
The relevant $\beta$ depends on sensor position and the instantaneous trajectory; it is not simply the vertical launch angle in every installation.
\item \textbf{Launch angles.}
Fit a trajectory only when the burst contains enough correctly timed and associated observations for the proposed state.
Filtering and smoothing can combine those observations, but cannot create information that the cadence or viewing geometry did not provide (\cref{sec:ekf}).
Compare the resulting launch angles against independent observations at a declared separation time, and test the benefit over a simpler estimator on held-out shots.
\item \textbf{Spin rate.}
Treat harmonic selection, surface-pattern ambiguity, observation duration and signal-to-noise ratio as explicit parts of the estimator (\cref{sec:radarspin}).
No universal 50--60\% detection ceiling follows from radar physics.
Longer windows can improve spectral resolution while also admitting changing geometry, acceleration and additional targets; comb or cepstral methods need tests for missed fundamentals and integer-multiple aliases.
A fallback from club, speed or loft priors is an impact-model prediction and must remain labeled as such.
Report its validation separately from successful rotation-sensitive observations.
\emph{Disclosure reference:} US8845442 describes radar spin processing; the technical map in \cref{sec:fto} does not determine clearance or a technique-wide expiry date.
\item \textbf{Spin axis.}
Short flight segments may provide too little trajectory information to separate axis, aerodynamic coefficients and launch-condition errors reliably.
Rotation-related observations from receiver geometry provide another disclosed route (\cref{sec:patentflightscope}), but require suitable hardware and validated signal processing.
Optical rotation estimates and explicitly labeled impact-model estimates are further options; their observation requirements and error sources differ.
Retain competing rotation hypotheses or withhold an axis when the observations do not resolve them.
\end{enumerate}

\section{Club Data: Define the Point, Frame and Time}

\begin{enumerate}
\item \textbf{Club speed.}
Define the reference point and event time before processing a pre-impact spectrum (\cref{sec:pathreference}).
For a rigid head, velocities at two points satisfy $\vect{v}_B=\vect{v}_A+\vect{\omega}\times(\vect{r}_B-\vect{r}_A)$.
A fixed spectral percentile may be a repeatable signal feature, but it does not by itself identify either point or remove the rotational contribution.
Reflectivity, viewing direction and shaft returns can change the percentile without the same change in the declared head-point speed.
Validate the mapping against tracked head motion over varied clubs and strikes.
Smash factor is a useful diagnostic only with compatible ball and club definitions; an idealized impact-model bound is not a universal hard rejection threshold (\cref{sec:smash}).
\item \textbf{Club path and attack angle.}
These describe the direction of velocity at the chosen point and time, derived from sufficient trajectory information in a calibrated frame.
A scattering-centre angle is not itself either quantity.
Alignment, association, differentiation, reference-point motion and timing all contribute uncertainty.
Check them separately rather than assuming alignment always dominates or that adding a radar makes path and attack angle direct observations.
\item \textbf{Face angle.}
The conditional D-plane inverse in \cref{eq:faceinversion} is a possible estimator, not a verified description of every commercial algorithm.
For a planar model $\ell=wf+(1-w)p$ with nonzero weight $w$, launch direction $\ell$, face angle $f$ and path $p$, first-order propagation gives
\begin{equation}
\label{eq:designfaceuncertainty}
\delta f\approx\frac{\delta\ell}{w}
-\frac{1-w}{w}\delta p
-\frac{\ell-p}{w^2}\delta w.
\end{equation}
The launch-to-face sensitivity $1/w$ is approximately 1.15 for $w=0.87$ and 1.32 for $w=0.76$; neither is a universal amplification.
Include covariance, weight uncertainty and model discrepancy, especially when strike location and gear effect violate the centered-impact approximation (\cref{sec:facepathweighting}).
If those conditions are unknown, enlarging uncertainty or withholding the result is more defensible than presenting the inverse as a face observation.
A shared forward/inverse software module still requires independent validation; algebraic consistency alone is insufficient.
\item \textbf{Impact location and orientation.}
Club and ball imaging offers a concrete route to locating contact relative to the face.
The Acushnet and Wintriss disclosures (\cref{sec:patentacushnet,sec:wintriss}) illustrate measurement architectures, rather than complete validated implementations or blanket rights clearance.
Face fiducials can constrain orientation through calibrated image geometry and a known marker-to-face relationship.
Camera count alone does not establish three-dimensional observability, visibility at impact or achieved accuracy.
Account for lens distortion, exposure blur, strobe timing, marker placement and the relation between the reconstructed surface and the face normal being reported.
\end{enumerate}

\section{Capability Tiers and Acceptance Conditions}

Each stage adds a testable capability to the preceding system.
An unsuccessful estimate at one stage must not be silently replaced by a more restrictive prior and then relabeled as a new measurement.

\begin{enumerate}
\item \textbf{Tier 1 --- validated radar observations.}
Establish calibrated radial velocity and whatever directional information the selected hardware actually provides.
Synchronize the trigger and observation windows, define ball/club association and preserve raw or sufficiently detailed intermediate data for replay.
Add spin or impact-model outputs only with provenance and separately tested failure criteria.
Acceptance requires a declared operating range, paired comparison results and per-parameter availability, rather than merely reproducing a consumer monitor's list of display fields.
\item \textbf{Tier 2 --- optical rotation and impact constraints.}
Add synchronized imaging where it addresses an identified ambiguity or missing observable.
Calibrated image sequences can constrain launch direction and surface rotation; impact location also requires an adequate view of the ball and club near contact.
A global shutter avoids row-to-row exposure timing differences but does not eliminate motion blur; illumination and exposure duration remain design variables.
Radar can contribute speed and triggering, while optical observations can distinguish rotation hypotheses.
Acceptance includes occlusion, marker symmetry, timing and calibration tests across the intended shot range, with costs measured for the actual implementation.
\item \textbf{Tier 3 --- validated joint estimation and prediction.}
Fuse complementary observations with sensor-specific measurement functions, timestamps and covariance models.
Shared calibration or processing can correlate their errors; counting the same information twice produces unjustifiably narrow uncertainty.
Version the flight model and propagate launch, aerodynamic and environmental uncertainty to the landing event (\cref{ch:flight}).
Acceptance requires checks that fusion improves held-out performance and uncertainty calibration, including prediction error on independently observed trajectories.
There is no universal yardage equivalent for a fixed coefficient change.
\item \textbf{Tier 4 --- observable club motion.}
Recover the club's pose and velocity from an observation geometry that supports the requested rigid-body state, then evaluate speed, path and attack angle at the declared point and time.
The screw-theoretic formulation in \cref{app:screw} provides the kinematic framework, but its rank and ambiguity analysis is part of the design requirement.
For a single monostatic radar phase centre, resolved Doppler rows alone do not independently determine all six instantaneous twist components; additional viewpoints, temporal position information or other observations are needed under stated motion assumptions.
Custom chirps and an $SE(3)$ smoother do not remove that ambiguity by themselves.
An instantaneous screw axis also becomes poorly conditioned near zero angular speed and is not automatically a global swing plane.
Optical face and impact measurements can be combined with the recovered motion, provided their frames and event times are aligned.
Acceptance requires reference-point, pose, twist and derived-axis validation, including the configurations in which the output must be withheld.
\end{enumerate}

\section{Reporting Honesty as a Design Feature}

A reported quantity should carry its definition, provenance, uncertainty and relevant model version.
The hierarchy in \cref{tab:hierarchy} organizes dependence on observations and assumptions; it is not an accuracy ranking.
Both a camera pose and a radar trajectory are estimates from sensor observations, and each can be wrong for different reasons.
The independent evidence in \cref{ch:accuracy} shows why return rate and agreement must be reported together, and why selected-shot precision cannot stand in for accuracy on all attempts.

A useful shot record distinguishes observed channels, model-conditioned estimates, simulated outcomes and prior-based fallbacks.
Preserve the coordinate frame, reference point, event time, calibration identifier, ball configuration and rejection reason alongside the value.
If an uncertainty interval is conditional on a fixed aerodynamic or impact model, state that condition.
Do not manufacture an uncertainty number where its calibration has not been tested.

These distinctions help the golfer interpret changes in the swing.
A change in predicted carry can originate in ball launch, atmosphere or software; a change in club speed can originate in a different reference point; an apparent face change can originate in the assumptions of an impact inverse.
The design objective is to make those alternatives testable with synchronized observations and a documented chain of inference.
That is what turns a plausible display into an instrument that can support research and coaching decisions.
